\chapter{Tự Học, Tự Quản Lý Và Trưởng Thành}

\begin{leadbox}
Từ một học sinh còn phụ thuộc đến một người có thể tự học là cả một hành trình dài.
Chương này đi vào những năng lực nhỏ nhưng tạo ra khác biệt rất lớn: biết lập kế hoạch, giữ lời hứa với chính mình, và đủ trung thực để tự kiểm điểm.
\end{leadbox}

\storycard{81. Tờ Kế Hoạch Đầu Tiên}{The First Study Plan}{Một kế hoạch được viết ra khiến tiến bộ trở nên nhìn thấy được, thay vì chỉ là một cảm giác mơ hồ.}
\storycard{82. Một Giờ Không Điện Thoại}{One Hour Without a Phone}{Tự chủ hiếm khi bắt đầu từ lời thề lớn; nó bắt đầu từ một khoảng thời gian nhỏ được giữ trọn vẹn.}
\storycard{83. Học Nhóm Mà Vẫn Hiểu Bài}{Studying Together Properly}{Học nhóm đúng nghĩa là cùng nâng nhau lên, không phải cùng kéo nhau ra khỏi sự tập trung.}
\storycard{84. Sổ Tay Sai Lầm}{The Mistake Notebook}{Ghi lại lỗi sai là cách biến thất bại thành tài liệu học tập quý giá nhất của chính mình.}
\storycard{85. Dậy Sớm Hơn 15 Phút}{Fifteen Minutes Earlier}{Kỷ luật cá nhân nhiều khi được dựng lên từ một thay đổi rất nhỏ nhưng được lặp lại đều đặn.}
\storycard{86. Đặt Mục Tiêu Thực Tế}{A Realistic Goal}{Mục tiêu đúng phải mở ra một con đường đi tới, chứ không dựng lên một bức tường khiến người học nản chí.}
\storycard{87. Tự Đánh Giá Sau Mỗi Tuần}{Weekly Self-Review}{Người biết nhìn lại mình sẽ bớt lệ thuộc vào việc luôn phải chờ người khác nhắc nhở.}
\storycard{88. Tìm Một Bạn Đồng Hành}{An Accountability Partner}{Động lực bền hơn khi có một người cùng nhắc nhau tiến bộ bằng sự tử tế và đều đặn.}
\storycard{89. Học Vì Ý Nghĩa}{Studying for Meaning}{Khi học sinh biết mình học để làm gì, sức bền sẽ mạnh hơn nhiều so với học chỉ để đối phó.}
\storycard{90. Tự Quản Lý Là Tự Do}{Self-Management Is Freedom}{Kỷ luật tự thân không phải xiềng xích; đó là cách một người mở rộng lựa chọn cho tương lai của mình.}
